\documentclass[11pt]{article}

\usepackage[a4paper,margin=2.5cm]{geometry}
\usepackage{amsmath,amssymb}
\usepackage{booktabs}
\usepackage{hyperref}
\usepackage{enumitem}
\usepackage{xcolor}
\usepackage{listings}

\lstset{
  basicstyle=\ttfamily\small,
  breaklines=true,
  columns=fullflexible,
  showstringspaces=false,
  frame=single,
}

\title{Coursework Report: PPO with Reward Shaping in RoboCasa Manipulation Tasks}
\date{April 27, 2026}

\begin{document}
\maketitle

\begin{abstract}
This report studies Proximal Policy Optimization (PPO) with reward shaping for two
robotic manipulation tasks in RoboCasa: \emph{apple-to-cabinet} and \emph{apple-to-bowl}.
The goal is to analyze how dense shaping terms affect learning progress and policy
behavior, especially in long-horizon sparse-reward settings. Experiments are run with
Stable-Baselines3 PPO and task-specific wrappers that expose success-related sub-metrics.
Results show that reward shaping is necessary to avoid complete learning collapse, but
also introduces local optima where the robot learns to grasp and hover near target
regions without consistently finishing placement and release. Section~\ref{sec:branch}
documents the changes introduced on branch
\texttt{feature/test-new-reward-shaping}, which target precisely this
``grasp-and-hover'' failure mode by adding an explicit \emph{hold} reward, a
\emph{close-while-near} ramp, an \emph{anti-open-under-contact} penalty, and strict
state gates configurable through a curriculum.
\end{abstract}

\section{Introduction}
Sparse-reward robotic manipulation remains challenging because meaningful learning
signals are rare and delayed. In long-horizon pick-and-place tasks, PPO may learn
slowly or collapse to ineffective exploration when only terminal success is rewarded.
Reward shaping addresses this issue by introducing intermediate objectives aligned
with task progress.

This report studies PPO with reward shaping in two RoboCasa tasks and evaluates how
shaping design affects (i) subskill acquisition and (ii) end-to-end success.

\section{Problem Formulation}

\subsection{Tasks}
We consider two related manipulation tasks:
\begin{itemize}
  \item \textbf{Apple-to-Cabinet (apple-to-cab):} pick an apple from a counter and
    place it into a cabinet region.
  \item \textbf{Apple-to-Bowl (apple-to-bowl):} pick an apple from a counter and
    place it into a bowl.
\end{itemize}

\subsection{Success Criteria and Transfer Challenge}
The two tasks share a pick--transport--place structure, but differ in terminal
constraints. The bowl task requires both successful placement and gripper
disengagement, making release behavior critical. As a result, policies transferred
from apple-to-cab may preserve grasp and transport skills yet still underperform on
final success in apple-to-bowl.

\section{Method}

\subsection{PPO Objective}
We use PPO with an MLP policy and clipped policy updates. PPO optimizes:
\begin{equation*}
  \mathcal{L}^{\text{CLIP}}(\theta) = \mathbb{E}_t \left[ \min\!\left(
  r_t(\theta)\hat{A}_t,\ \mathrm{clip}(r_t(\theta), 1-\epsilon, 1+\epsilon)\hat{A}_t
  \right) \right] .
\end{equation*}

\subsection{Reward Shaping Framework}
The shaped reward augments sparse environment reward with dense terms for:
\begin{itemize}
  \item contact and grasp acquisition,
  \item lift and transport progress,
  \item release and placement events,
  \item final success bonus.
\end{itemize}

For state labeling, v2/v3 do not directly use RoboCasa's default
\texttt{OU.check\_obj\_grasped}. Instead, they use a task-specific robust grasp
check and then define a stricter \emph{effective grasped} state through additional
gating (e.g., stability/transportability conditions). This redefinition is used
consistently by reward logic and episode metrics. The implementation also logs
atomic metrics such as \texttt{contact\_success\_rate},
\texttt{grasp\_success\_rate}, \texttt{grasp\_lift\_success\_rate}, and
\texttt{place\_success\_rate} to inspect stage-wise learning.

\subsection{Lightweight Reward Curriculum Across Versions}
The shaping pipeline uses a lightweight reward-curriculum mechanism across v1--v3.
Specifically, dense shaping is progressively scaled during early training (curriculum
steps), so optimization is not immediately dominated by shaped rewards. In addition,
phase-capping is available in later versions (v2/v3) to restrict active shaping
terms to earlier skill stages (e.g., lift-focused phase) before enabling the full
objective set. This curriculum is reward-centric (not environment-switching
curriculum), and is designed to stabilize early exploration while preserving the
final task objective.

\subsection{Versioned Reward Design and Term Interpretation}
Table~\ref{tab:rewards} compares core reward coefficients for the three report
versions. This comparison highlights the shift from moderate shaping (v1), to
grasp-stability shaping (v2), to strongly outcome-biased shaping for bowl transfer
(v3).

\begin{table}[h]
\centering
\caption{Key reward-term defaults across v1, v2, and v3.}
\label{tab:rewards}
\begin{tabular}{lccc}
\toprule
Reward Term & v1 & v2 & v3 \\
\midrule
Reach reward         & 1.5 & 3.0 & 3.0 \\
Contact reward       & N/A & 0.5 & 0.5 \\
Pre-grasp reward     & 0.1 & 0.3 & 0.3 \\
Grasp reward         & 3.0 & 3.0 & 3.0 \\
Grasp-hold reward    & N/A & 0.3 & 0.01 \\
Lift reward          & 2.0 & 2.0 & 0.5 \\
Lift-sustain reward  & N/A & 0.5 & 2.0 \\
Release reward       & 0.5 & 0.5 & 2.0 \\
Transport reward     & 2.0 & 2.0 & 6.0 \\
Place reward         & 3.0 & 3.0 & 3.0 \\
Success reward       & 10.0 & 10.0 & 100.0 \\
\bottomrule
\end{tabular}
\end{table}

In addition, v3 introduces explicit anti-cheating penalties and a retraction
reward, which are not present in v1/v2. This makes v3's objective more strongly
coupled to complete placement behavior rather than grasp-only progress. Below is a
concise interpretation of the main reward terms used across versions:

\begin{itemize}
  \item \textbf{Reach reward}: encourages the end-effector to move closer to the object.
  \item \textbf{Contact reward}: gives credit when the gripper first establishes object contact.
  \item \textbf{Pre-grasp reward}: rewards near-object alignment before a stable grasp is formed.
  \item \textbf{Grasp reward}: provides a bonus when a valid grasp is achieved.
  \item \textbf{Grasp-hold reward}: rewards maintaining grasp stability over time.
  \item \textbf{Lift reward}: encourages lifting the object above its initial height.
  \item \textbf{Lift-sustain reward}: rewards keeping the object elevated after lifting.
  \item \textbf{Release reward}: encourages opening the gripper at the target placement stage.
  \item \textbf{Transport reward}: rewards moving the grasped object toward the goal region.
  \item \textbf{Place reward}: gives credit when the object reaches the target receptacle.
  \item \textbf{Success reward}: sparse terminal bonus for satisfying full task completion criteria.
  \item \textbf{Drop / anti-cheating penalties}: penalize undesirable shortcuts (e.g., premature lift, air-close, accidental drop), improving policy robustness.
  \item \textbf{Retraction reward (v3)}: rewards pulling the gripper away after placement to satisfy strict success checks.
\end{itemize}

\section{Experimental Protocol}

\subsection{Software Stack}
\begin{itemize}
  \item Simulator stack: RoboCasa + robosuite.
  \item RL library: Stable-Baselines3 PPO.
  \item Training implementation: PPO reward-shaping v3 entrypoint.
  \item Evaluation implementation: deterministic policy rollout entrypoint.
\end{itemize}

\subsection{Training Configuration}
Representative settings used in the main reported v3 run (not necessarily script
defaults):
\begin{itemize}
  \item horizon: $700$,
  \item $n\_steps = 2048$, batch size $= 256$, $n\_envs = 4$,
  \item learning rate schedule from $5\times 10^{-5}$ to $5\times 10^{-6}$,
  \item v3 is initialized from a pretrained v2 checkpoint (policy and
    normalization state),
  \item checkpoint-based resume is supported, enabling breakpoint recovery and
    continued training,
  \item total timesteps up to $\sim 1.56\text{M}$ for v3 logs analyzed here.
\end{itemize}

\section{Results}

\subsection{Quantitative Comparison}
Table~\ref{tab:results} summarizes final logged performance for the three
reported model versions.

\begin{table}[h]
\centering
\caption{Final performance comparison across v1--v3 (from available logs in this repository).}
\label{tab:results}
\begin{tabular}{llrrrr}
\toprule
Version & Task & Timestep & Success Rate & Place Success & Grasp Success \\
\midrule
v1 & apple-to-cab  & 3{,}006{,}464 & 0\%  & 0\%  & 81\%  \\
v2 & apple-to-cab  & 3{,}006{,}464 & 0\%  & 0\%  & 100\% \\
v3 & apple-to-bowl & 1{,}556{,}480 & 8\%  & 13\% & 82\%  \\
\bottomrule
\end{tabular}
\end{table}

\subsection{Key Observations}
\begin{itemize}
  \item v1 and v2 both show strong grasp learning on apple-to-cab, with v2
    reaching near-saturated grasp rate.
  \item v3 achieves non-zero placement and final success on apple-to-bowl,
    unlike v1/v2 in their reported logs.
  \item The gap between grasp-related and success-related metrics in v3
    indicates a late-stage bottleneck (release, stabilization, or
    disengagement).
\end{itemize}

\subsection{Behavioral Evidence from Rollouts}
Observed policy behavior in evaluation videos matches the metric trends:
\begin{itemize}
  \item the gripper often reaches and secures the apple,
  \item the end-effector transports near the bowl region,
  \item but episodes frequently fail to complete a clean release-away sequence.
\end{itemize}
This suggests a local optimum: maximizing dense transport/grasp incentives without
reliably satisfying strict terminal conditions.

\section{Discussion}

\subsection{Why Shaping Helps}
Compared with sparse-only learning, shaped rewards clearly bootstrap useful
skills (contact, grasp, transport), making training dynamics non-trivial and
reproducible.

\subsection{Why Final Success Remains Limited}
Three likely factors:
\begin{enumerate}
  \item \textbf{Reward imbalance}: sustained intermediate rewards can dominate rare
    terminal rewards. In practice, dense shaping terms (e.g., reach, grasp,
    transport, and hold-related rewards) are collected frequently within an
    episode, so their cumulative contribution to return can become much larger
    than the sparse success bonus that appears only at full completion. Under
    this imbalance, the policy may converge to a locally optimal strategy that
    reliably maximizes repeatable partial-progress rewards (grasping and
    hovering near the target) without consistently executing the high-precision
    release-and-disengage sequence required for final success.
  \item \textbf{Strict success definition}: bowl placement also requires release
    and sufficient gripper--object separation.
  \item \textbf{Training horizon may be insufficient}: the current runs may
    require longer training steps to clearly observe whether final success
    continues to improve or plateaus.
\end{enumerate}

These factors likely interact rather than act in isolation. In particular, transfer
from a strong v2 grasp-centric prior can stabilize early subskills but also preserve
behavior patterns that are suboptimal for bowl completion (e.g., grasp-and-hover
tendencies). Combined with strict terminal checks and dense intermediate rewards,
this interaction can produce the observed metric profile: relatively strong grasp
rates but lower place/final success rates. \textbf{Despite introducing a lightweight
reward curriculum}, the current training dynamics still indicate that late-stage
release-and-disengagement behavior remains the principal bottleneck.

\subsection{Limitations}
This report focuses on logged training curves from a limited set of runs. No broad
hyperparameter sweep or significance testing is included; therefore conclusions
should be treated as empirical findings rather than definitive generalization.

\section{Conclusion}
This report shows that PPO with reward shaping can reliably learn manipulation
\emph{subskills} (especially reach, grasp, and transport) in RoboCasa tasks, but
full task completion remains substantially harder than intermediate progress.
Across the analyzed runs, the apple-to-bowl setting achieves non-zero success yet
still exhibits a clear gap between high grasp-related metrics and lower final
success metrics.

The results indicate that this gap is mainly associated with late-stage execution
requirements: successful bowl completion depends not only on placing the object in
the receptacle, but also on a clean release-and-disengagement sequence. Under dense
shaping, policies can therefore converge to local optima that favor repeatable
partial progress over precise terminal behavior.

Overall, the key lesson is that reward shaping is necessary but not sufficient for
robust end-to-end success in strict long-horizon manipulation tasks. Future work
should focus on improving post-placement reward balance and extending training
duration to better determine whether final success continues to improve or reaches
an early plateau.

% =====================================================================
% New content: branch feature/test-new-reward-shaping
% =====================================================================
\section{Branch \texttt{feature/test-new-reward-shaping}: Targeted Fixes for Grasp-and-Hold Instability}
\label{sec:branch}

This section documents the changes introduced on the branch
\texttt{feature/test-new-reward-shaping} (commits
\texttt{a6ad4e6}, \texttt{85e424b}, \texttt{a61e9a7}, \texttt{2646c74}). The
focus is to address the grasp-and-hold instability identified in the v3 rollouts
of Section~5.3 and to make the experimental protocol reproducible on a SLURM
cluster.

\subsection{Diagnostic: why grasping was unstable}
Evaluation logs of v3 showed that the policy reached the apple but rarely
transitioned into a stable \emph{grasp-and-hold} regime: \texttt{grasp} and
\texttt{inside} indicators were near zero across most episodes
(\texttt{eval\_robocasa\_dense.py}: \texttt{grasp=0.000}, \texttt{inside=0.000}
in most steps). The dominant failure mode was a release-then-drop pattern:
the gripper closed, then partially reopened during transport. Two contributing
mechanisms were hypothesized:
\begin{itemize}
  \item \textbf{Gripper-action jitter}: PPO exploration noise on the gripper
    dimension, combined with the absence of an explicit cost on opening under
    contact, produced micro-oscillations that broke the contact streak.
  \item \textbf{Insufficient hold incentive}: the v3 reward paid for
    \emph{achieving} a grasp but barely for \emph{maintaining} it
    (Table~\ref{tab:rewards}: grasp-hold weight $0.01$).
\end{itemize}

\subsection{New shaping wrapper: \texttt{DenseStrictRewardWrapper}}
A new training entrypoint
\texttt{scripts/train\_ppo\_reward\_shaping\_dense.py} introduces
\texttt{DenseStrictRewardWrapper}, a Gym wrapper that recomputes the dense
shaped reward each step and exposes per-term values through
\texttt{info["dense\_terms"]}. The same wrapper is used by
\texttt{scripts/eval\_robocasa\_dense.py} so that evaluation reward and
training reward are strictly comparable. The shaping vector is:
\begin{equation*}
r_{\text{dense}} = \underbrace{w_r e^{-\tau_r d_{eef,obj}}}_{\text{reach}}
+ \underbrace{w_g\, \mathbb{1}_{\text{grasped}}}_{\text{grasp}}
+ \underbrace{w_h\, \tfrac{s_{\text{grasp}}}{S_{\max}}}_{\text{hold}}
+ \underbrace{w_n \cdot \mathrm{near}(d) \cdot c_{\text{grip}}}_{\text{close-near}}
- \underbrace{w_o \cdot \mathbb{1}_{\text{contact}} \cdot \Delta_{\text{open}}}_{\text{open-contact penalty}}
+ \underbrace{w_l\, \tfrac{h}{h_{\min}}}_{\text{lift}}
+ \underbrace{w_c e^{-\tau_c d_{obj,\text{tgt}}}}_{\text{carry}}
+ \underbrace{w_i\, \mathbb{1}_{\text{inside}}}_{\text{inside}}
+ \underbrace{w_s\, \mathbb{1}_{\text{success}}}_{\text{success}}
\end{equation*}
where $s_{\text{grasp}}$ is the consecutive-grasp streak, $c_{\text{grip}}$ is the
gripper closure fraction in $[0,1]$, $\Delta_{\text{open}}$ is the per-step decrease
of $c_{\text{grip}}$, and $\mathrm{near}(d) = \mathrm{clip}(1 - d/d_{\text{near}}, 0, 1)$.

The four ideas listed in the project brief map to specific terms:
\begin{itemize}
  \item \emph{Hold reward (most important)} $\to$ \texttt{w\_hold} on the
    grasp streak \texttt{raw\_grasp\_streak / hold\_max\_steps}.
  \item \emph{Close-while-near} $\to$ \texttt{w\_close\_near}, active even
    \emph{before} \texttt{check\_obj\_grasped} flips, so the gradient field
    pulls the policy toward closing as it approaches the object.
  \item \emph{Open-under-contact penalty} $\to$ \texttt{w\_open\_contact\_penalty}
    weighting $\Delta_{\text{open}}$ whenever
    \texttt{raw\_grasped or raw\_grasp\_streak > 0}.
  \item \emph{Strict gates softened at start, then hardened} $\to$
    \texttt{StrictStateCfg} parameters
    (\texttt{grasp\_min\_consecutive}, \texttt{grasp\_min\_close},
    \texttt{inside\_min\_consecutive}, \dots) made fully configurable through
    CLI flags so a curriculum can ramp them.
\end{itemize}

\subsubsection*{New shaping hyperparameters and their defaults}
Table~\ref{tab:new-shaping} lists the new shaping weights and their script
defaults (\texttt{train\_ppo\_reward\_shaping\_dense.py},
\texttt{train\_ppo\_curriculum\_dense.py}). These weights are intentionally
small relative to v3 success bonus (\texttt{w\_success}=5.0) so that the dense
terms guide exploration without dominating the terminal signal.

\begin{table}[h]
\centering
\caption{New dense-reward weights introduced on the branch.}
\label{tab:new-shaping}
\begin{tabular}{lll}
\toprule
Flag & Default & Role \\
\midrule
\texttt{--w\_reach}                 & 0.25 & exponential reach shaping \\
\texttt{--w\_grasp}                 & 0.35 & strict-grasped indicator \\
\texttt{--w\_hold}                  & 0.25 & grasp-streak hold bonus \\
\texttt{--w\_close\_near}           & 0.15 & close-while-near ramp \\
\texttt{--w\_open\_contact\_penalty}& 0.25 & negative for opening under contact \\
\texttt{--w\_lift}                  & 0.20 & normalized lift height \\
\texttt{--w\_carry}                 & 0.35 & exponential carry to target \\
\texttt{--w\_inside}                & 0.75 & strict-inside indicator \\
\texttt{--w\_success}               & 5.0  & terminal success bonus \\
\texttt{--reach\_temp}              & 4.0  & decay rate of reach shaping \\
\texttt{--carry\_temp}              & 4.0  & decay rate of carry shaping \\
\texttt{--close\_near\_dist\_m}     & 0.05 & ramp width for close-while-near (m) \\
\texttt{--hold\_max\_steps}         & 50   & saturation horizon of hold streak \\
\bottomrule
\end{tabular}
\end{table}

\begin{table}[h]
\centering
\caption{Strict state-gate parameters (\texttt{StrictStateCfg}).}
\label{tab:strict}
\begin{tabular}{lll}
\toprule
Flag & Default & Role \\
\midrule
\texttt{--strict\_grasp\_lift\_m}         & 0.015 & min lift to count as transportable (m) \\
\texttt{--strict\_grasp\_min\_consecutive}& 4     & consecutive raw-grasped steps \\
\texttt{--strict\_grasp\_follow\_dist\_m} & 0.035 & max eef-obj distance during follow (m) \\
\texttt{--strict\_gripper\_open\_qpos}    & 0.040 & qpos used to normalize closure \\
\texttt{--strict\_grasp\_min\_close}      & 0.60  & min closure fraction \\
\texttt{--strict\_inside\_min\_consecutive}& 6    & inside streak required \\
\texttt{--strict\_inside\_max\_obj\_speed}& 0.15  & max obj speed while ``inside'' (m/s) \\
\texttt{--allow\_inside\_while\_grasped}  & off   & if set, success allowed while still grasped \\
\bottomrule
\end{tabular}
\end{table}

\subsection{Curriculum entrypoint}
A second entrypoint, \texttt{train\_ppo\_curriculum\_dense.py}, wraps the
environment with the existing \texttt{CurriculumWrapper} \emph{before} the
\texttt{DenseStrictRewardWrapper} and adds a \texttt{CurriculumCallback} that
advances stages on a rolling success-rate criterion:
\begin{itemize}
  \item \texttt{--curriculum\_window}=100 episodes,
  \item \texttt{--curriculum\_min\_timesteps}=50{,}000 per stage,
  \item \texttt{--curriculum\_thresholds}=``0.70,0.80'' (advance when window
    success exceeds threshold).
\end{itemize}
This expresses the brief's idea of \emph{starting easy, then hardening}: stage
$0$ exposes the easier subgoal, and the strict gates (which the operator can
ramp via CLI) become operative once the policy is competent.

\subsection{New PPO knobs (curriculum entrypoint)}
On top of the v3 hyperparameters, the curriculum script exposes additional
PPO controls intended to reduce gripper jitter and improve stability:
\begin{itemize}
  \item \texttt{--n\_epochs} (default 10), \texttt{--gamma} (0.995),
    \texttt{--gae\_lambda} (0.95),
  \item \texttt{--ent\_coef}, \texttt{--clip\_range} (0.2),
    \texttt{--vf\_coef} (0.5), \texttt{--max\_grad\_norm} (0.5),
  \item \texttt{--target\_kl} (optional early-stop on KL divergence).
\end{itemize}
Default \texttt{--learning\_rate} is $3\times 10^{-4}$,
\texttt{--n\_steps} $= 2048$, \texttt{--batch\_size} $= 256$, matching the
report v3 baseline so changes can be ablated cleanly.

\subsection{Evaluation script alignment}
\texttt{eval\_robocasa\_dense.py} reuses \texttt{DenseStrictRewardWrapper} and
\texttt{StrictStateCfg} so that ``what the policy sees at training time'' and
``what we measure at eval'' are identical wrappers. It also:
\begin{itemize}
  \item renders a $2 \times 2$ camera tile
    (\texttt{robot0\_agentview\_center/left/right} +
    \texttt{robot0\_eye\_in\_hand}) for diagnosis,
  \item prints per-episode \texttt{dense\_terms} so that breakdowns (reach,
    grasp, hold, close-near, open-contact, lift, carry, inside, success) are
    visible per rollout,
  \item exposes the same strict-gate flags as training to support sensitivity
    analysis without retraining.
\end{itemize}

\subsection{SLURM Experimental Protocol}
\label{sec:slurm}

\subsubsection*{Cluster topology and remote layout}
Jobs run on the partition \texttt{3090} (single RTX 3090 per job). The remote
working tree is \texttt{/home/infres/latoundji-25/robocasa/RoboCasa-RL}.
Synchronisation is one-way (local $\to$ remote) over rsync via the helper
scripts in \texttt{scripts/}:
\begin{lstlisting}
./scripts/cluster_sync_push.sh        # sources -> cluster (excludes .venv, runs/, artifacts/)
./scripts/cluster_exec.sh submit train.sh --export=KEY=VAL,...
./scripts/cluster_exec.sh status      # squeue --me
./scripts/cluster_exec.sh logs <jobid>
./scripts/cluster_sync_pull.sh        # cluster -> local (runs/, models/, eval_videos/)
\end{lstlisting}

\subsubsection*{SLURM resource request (\texttt{train.sh})}
\begin{itemize}
  \item \texttt{--partition=3090}, \texttt{--gres=gpu:1},
  \item \texttt{--mem=32G} (raised from the initial 8\,GB after observing OOM
    kills with $n\_envs{=}8$ \texttt{SubprocVecEnv}; each worker loads
    MuJoCo+robocasa+torch in its own address space),
  \item \texttt{--cpus-per-task=9} (1 main + 8 workers),
  \item \texttt{--time=24:00:00}, \texttt{--output=runs/slurm/\%x\_\%j.out}.
\end{itemize}

\subsubsection*{Environment bootstrap}
\texttt{train.sh} bootstraps a persistent uv-managed venv at
\texttt{\$HOME/.venvs/robocasa\_rl\_<partition>}, separate from the repo
\texttt{.venv}, so concurrent jobs share a cache. The bootstrap is serialised
through an \texttt{flock} on \texttt{.uv\_sync.lock} and consists of:
\begin{enumerate}
  \item \texttt{uv venv --python python3.11} (first run only),
  \item \texttt{uv sync --frozen --no-dev --inexact} (lockfile-driven, keeps
    extras across runs),
  \item \texttt{uv pip install --python \$VENV/bin/python <extras>} where the
    extras list mirrors the working local venv:
    \texttt{termcolor, mujoco==3.3.1, numpy==2.2.5, scipy, opencv-python,
    pyyaml, pillow, pygame, qpsolvers[quadprog], pynput, tqdm, rich, h5py,
    lxml, tensorboard, imageio, imageio-ffmpeg, av, matplotlib, psutil},
  \item idempotent copy
    \texttt{robosuite/macros.py} $\to$ \texttt{robosuite/macros\_private.py}
    (replaces the upstream interactive setup script, unusable in batch).
\end{enumerate}
The hard pins \texttt{mujoco==3.3.1} and \texttt{numpy==2.2.5} are required by
\texttt{robocasa/\_\_init\_\_.py} \texttt{assert} statements; do \emph{not}
add \texttt{mink} or \texttt{robosuite\_models} as they pull \texttt{numpy<2}
and break the assert.

\subsubsection*{Reproducible launch command}
The reference run for the branch is:
\begin{lstlisting}
./scripts/cluster_sync_push.sh
./scripts/cluster_exec.sh submit train.sh \
  --export=PYTHON_BIN=python3.11,RUN_NAME=reward_dense_${USER},\
TOTAL_TIMESTEPS=10000000,N_ENVS=8,DEVICE=cpu
\end{lstlisting}
\texttt{DEVICE=cpu} reflects the SB3 advisory that PPO with an
\texttt{MlpPolicy} is typically faster on CPU than on GPU; this is verified
empirically per partition.

\subsubsection*{Logging and artifacts}
Each run writes to \texttt{models/\$RUN\_NAME/} on the cluster:
\begin{itemize}
  \item \texttt{config.json}: every CLI argument used,
  \item \texttt{logs/tensorboard/}: SB3 + \texttt{debug/*} metrics
    (\texttt{--log\_debug}),
  \item \texttt{logs/monitor/env\_<rank>/}: per-env Monitor CSVs,
  \item \texttt{logs/metrics.csv}: rollout-level summary
    (\texttt{ep\_rew\_mean}, \texttt{ep\_len\_mean},
    \texttt{success\_rate}),
  \item \texttt{checkpoints/} every \texttt{--checkpoint\_freq} steps,
  \item \texttt{best/best\_model.zip} every \texttt{--best\_eval\_freq} steps
    (greedy 5-episode eval),
  \item \texttt{ppo\_final.zip} at the end of training.
\end{itemize}

\subsection{Summary of changes by file}
\begin{itemize}
  \item \texttt{scripts/train\_ppo\_reward\_shaping\_dense.py} (new):
    \texttt{DenseStrictRewardWrapper}, \texttt{StrictStateCfg},
    \texttt{MetricsLoggerCallback}, \texttt{RolloutDebugLoggerCallback}.
  \item \texttt{scripts/train\_ppo\_curriculum\_dense.py} (new):
    \texttt{CurriculumCallback} on top of the dense wrapper, additional PPO
    knobs.
  \item \texttt{scripts/eval\_robocasa\_dense.py} (new): training-aligned
    evaluation, multi-camera tiled video, per-term reward dump.
  \item \texttt{train.sh} / \texttt{eval.sh} (new): SLURM headers, persistent
    uv venv bootstrap, macro-private setup, sanity checks for
    \texttt{termcolor / mujoco / numpy} versions.
  \item \texttt{scripts/cluster\_exec.sh}, \texttt{cluster\_sync\_push.sh},
    \texttt{cluster\_sync\_pull.sh} (new): submit/status/logs/cancel
    wrappers and rsync helpers.
  \item \texttt{\_common.sh} (new): shared SLURM environment setup.
  \item \texttt{env/custom\_pnp\_counter\_to\_cab.py},
    \texttt{scripts/train\_ppo\_curriculum.py},
    \texttt{docs/ALGORITHM\_FLOW.md}: minor edits supporting the new
    curriculum entrypoint.
\end{itemize}

\subsection{Expected impact and what to look for in the logs}
The new shaping is designed so that, in TensorBoard
(\texttt{--log\_debug}), one should observe:
\begin{itemize}
  \item rising \texttt{debug/dense\_terms/hold} as the policy learns to keep
    the gripper closed,
  \item near-zero \texttt{debug/dense\_terms/open\_contact\_penalty} once
    the policy stops oscillating on the gripper dimension,
  \item rising \texttt{debug/raw\_grasp\_streak} (longer holds) before
    \texttt{debug/dense\_terms/inside} starts firing,
  \item \texttt{curriculum/stage} stepping up only after
    \texttt{curriculum/success\_rate\_window} crosses the configured
    thresholds.
\end{itemize}
Failure to observe these patterns under the default weights is itself a
useful signal: it indicates that the strict gates are still too tight at the
current stage and should be relaxed via CLI before re-running.

\section*{Reproducibility Notes}
\begin{itemize}
  \item Main components: PPO reward-shaping v3 training pipeline, dense+strict
    wrapper, curriculum entrypoint, deterministic evaluation pipeline.
  \item Data sources: versioned training metrics logs for the v1, v2, and v3
    experiments; CSVs and TensorBoard runs under \texttt{models/<run\_name>/logs}.
  \item Evaluation media: rollout videos generated from deterministic policy
    inference, $2 \times 2$ tiled cameras, saved under
    \texttt{eval\_videos/<run\_name>/}.
  \item Branch: \texttt{feature/test-new-reward-shaping}, commits
    \texttt{a6ad4e6} \dots \texttt{2646c74}.
\end{itemize}

\end{document}
